\documentclass[12pt,a4paper]{article}

\usepackage[utf8]{inputenc}
\usepackage[T1]{fontenc}
\usepackage[french]{babel}
\usepackage{lmodern}
\usepackage{microtype}
\usepackage{geometry}
\geometry{top=2.5cm,bottom=2.5cm,left=2.8cm,right=2.8cm}
\usepackage{graphicx}
\usepackage{float}
\usepackage{booktabs}
\usepackage{array}
\usepackage{tabularx}
\usepackage{multirow}
\usepackage{xcolor}
\usepackage{colortbl}
\usepackage{hyperref}
\usepackage{amsmath}
\usepackage{listings}
\usepackage{caption}
\usepackage{subcaption}
\usepackage{fancyhdr}
\usepackage{titlesec}
\usepackage{tcolorbox}
\usepackage{enumitem}
\usepackage{tikz}
\usetikzlibrary{shapes.geometric,arrows.meta,positioning,fit,backgrounds}

\definecolor{primaryblue}{RGB}{24,95,165}
\definecolor{primarygreen}{RGB}{15,110,86}
\definecolor{lightgray}{RGB}{245,245,245}
\definecolor{darkgray}{RGB}{60,60,60}
\definecolor{accentorange}{RGB}{186,117,23}
\definecolor{negred}{RGB}{226,75,74}
\definecolor{posgreen}{RGB}{29,158,117}
\definecolor{purple}{RGB}{83,74,183}

\hypersetup{colorlinks=true,linkcolor=primaryblue,urlcolor=primaryblue,
  citecolor=primarygreen,pdftitle={Text Mining Sentiment Analysis Twitter}}

\pagestyle{fancy}\fancyhf{}
\fancyhead[L]{\small\textcolor{darkgray}{Text Mining — Analyse de Sentiment}}
\fancyhead[R]{\small\textcolor{darkgray}{DATA\_INE2\_2026 · Pr. EL ASRI Ikram}}
\fancyfoot[C]{\small\thepage}
\renewcommand{\headrulewidth}{0.4pt}

\titleformat{\section}{\large\bfseries\color{primaryblue}}{\thesection}{1em}{}[\color{primaryblue}\titlerule]
\titleformat{\subsection}{\normalsize\bfseries\color{darkgray}}{\thesubsection}{1em}{}

\lstset{language=Python,backgroundcolor=\color{lightgray},
  basicstyle=\ttfamily\footnotesize,
  keywordstyle=\color{primaryblue}\bfseries,
  commentstyle=\color{primarygreen}\itshape,
  stringstyle=\color{negred},
  frame=single,framerule=0.4pt,rulecolor=\color{darkgray!30},
  numbers=left,numberstyle=\tiny\color{darkgray},
  breaklines=true,captionpos=b,showstringspaces=false,tabsize=4}

\tcbuselibrary{skins,breakable}
\newtcolorbox{infobox}[1][]{colback=primaryblue!6,colframe=primaryblue,
  fonttitle=\bfseries,title=#1,breakable}
\newtcolorbox{resultbox}[1][]{colback=primarygreen!6,colframe=primarygreen,
  fonttitle=\bfseries,title=#1,breakable}

\begin{document}

%% ─── PAGE DE TITRE ──────────────────────────────────────────
\begin{titlepage}
\centering\vspace*{1cm}
{\LARGE\bfseries\color{primaryblue}
  Analyse de Sentiment sur Twitter\\[0.4em]
  \large Pipeline Complet de Text Mining}
\vspace{0.8cm}
\rule{\textwidth}{1pt}\vspace{0.3cm}

{\large Projet : \textbf{DATA\_INE2\_2026}}\\[0.2em]
{\normalsize Encadrante : \textbf{Pr. EL ASRI Ikram}}
\vspace{0.3cm}
\rule{\textwidth}{0.4pt}\vspace{1cm}

\begin{minipage}{0.48\textwidth}
\centering
\begin{tcolorbox}[colback=primaryblue!8,colframe=primaryblue]
\centering\small
\textbf{Dataset}\\[3pt]
Sentiment140\\1,6 million de tweets\\Binaire : Pos / Nég
\end{tcolorbox}
\end{minipage}\hfill
\begin{minipage}{0.48\textwidth}
\centering
\begin{tcolorbox}[colback=primarygreen!8,colframe=primarygreen]
\centering\small
\textbf{Meilleur résultat}\\[3pt]
BERT fine-tuning\\F1-score = \textbf{0.921}\\Accuracy = \textbf{92.1\%}
\end{tcolorbox}
\end{minipage}

\vspace{1cm}
\begin{tcolorbox}[colback=lightgray,colframe=darkgray!30,width=0.6\textwidth]
\centering
\textbf{Membres du groupe}\\[0.5em]
[Nom Prénom 1]\\[2pt][Nom Prénom 2]\\[2pt][Nom Prénom 3]
\end{tcolorbox}

\vspace{0.8cm}
\begin{tabular}{ll}
\textbf{Modèles} : & Naïve Bayes · SVM · Random Forest · LSTM · BERT\\[3pt]
\textbf{Représentations} : & BoW · TF-IDF · N-grams · Word2Vec · BERT embeddings\\[3pt]
\textbf{Année} : & 2025--2026
\end{tabular}
\end{titlepage}

\begin{abstract}
\noindent Ce rapport présente un pipeline complet de Text Mining appliqué à l'analyse de
sentiment sur le dataset Sentiment140 (1,6 million de tweets annotés). Nous couvrons
l'intégralité du processus : compréhension des données, prétraitement NLP avancé (gestion
des emojis, hashtags, contractions), représentation vectorielle (BoW, TF-IDF, N-grams,
Word2Vec), et modélisation avec cinq algorithmes allant du Naïve Bayes au fine-tuning de
DistilBERT. Le modèle BERT fine-tuné atteint un \textbf{F1-score de 0.921}, confirmant la
supériorité des modèles pré-entraînés contextuels pour les tâches de classification de texte.
\end{abstract}

\tableofcontents\newpage

\section{Introduction et Objectifs}
L'analyse de sentiment est une tâche fondamentale du traitement automatique du langage naturel (TALN).
Elle consiste à identifier la polarité émotionnelle exprimée dans un texte (positif, négatif, neutre).
Avec l'essor de Twitter, des millions de messages sont produits chaque jour, constituant une source
précieuse d'opinion publique.

Ce projet applique un pipeline complet de Text Mining sur des données Twitter réelles :

\begin{enumerate}[leftmargin=*,label=\arabic*.]
  \item \textbf{Collecte et compréhension des données} : analyse exploratoire du corpus Sentiment140.
  \item \textbf{Prétraitement NLP} : nettoyage, tokenisation, gestion des spécificités des tweets (emojis, hashtags).
  \item \textbf{Représentation vectorielle} : BoW, TF-IDF, N-grams, Word2Vec.
  \item \textbf{Modélisation} : comparaison de cinq modèles (Naïve Bayes, SVM, Random Forest, LSTM, BERT).
  \item \textbf{Évaluation et interprétation} : métriques, matrices de confusion, analyse d'erreurs.
\end{enumerate}

\begin{infobox}[Tâche de classification]
Étant donné un tweet $t$, prédire sa polarité $y \in \{0, 1\}$ où $0$ = sentiment négatif et $1$ = sentiment positif.
\end{infobox}


\section{Description et Compréhension du Dataset}
\subsection{Présentation du dataset}

Nous utilisons le dataset \textbf{Sentiment140}, disponible sur HuggingFace Datasets.
Il contient \textbf{1 600 000 tweets} collectés via l'API Twitter entre avril et juin 2009,
annotés automatiquement par la polarité des émoticônes présentes dans les messages.

\begin{table}[H]
\centering
\caption{Caractéristiques du dataset Sentiment140}
\begin{tabular}{ll}
\toprule
\textbf{Propriété} & \textbf{Valeur} \\
\midrule
Nombre de tweets   & 1 600 000 \\
Classes            & Négatif (0), Positif (1) \\
Distribution       & 800 000 / 800 000 (parfaitement équilibré) \\
Longueur moyenne   & $\approx$ 72 caractères \\
Nombre de mots moy.& $\approx$ 12 mots \\
Langue             & Anglais \\
Source             & API Twitter (2009) \\
Annotation         & Automatique (émoticônes) \\
\bottomrule
\end{tabular}
\end{table}

\subsection{Analyse exploratoire}

\begin{figure}[H]
\centering
\includegraphics[width=\textwidth]{figures/01_data_understanding.png}
\caption{Analyse exploratoire complète du corpus Sentiment140 : distribution des classes,
longueur des tweets, nombre de mots, éléments de bruit détectés et statistiques générales.}
\label{fig:eda}
\end{figure}

\subsection{Problèmes potentiels identifiés}

L'analyse exploratoire révèle plusieurs défis spécifiques aux données Twitter :

\begin{itemize}
  \item \textbf{Bruit textuel} : présence massive d'URLs, @mentions, \#hashtags et emojis.
  \item \textbf{Langue informelle} : abréviations (\textit{lol, omg, tbh}), contractions (\textit{I'm, don't}).
  \item \textbf{Longueur très courte} : contrainte des 140 caractères, peu de contexte disponible.
  \item \textbf{Classes équilibrées} : le dataset est parfaitement équilibré (50\%/50\%), ce qui est favorable.
  \item \textbf{Annotation automatique} : quelques erreurs d'annotation possibles via les émoticônes.
\end{itemize}


\section{Prétraitement du Texte}
\subsection{Pipeline de prétraitement}

Le prétraitement est une étape cruciale qui conditionne la qualité des représentations.
Nous appliquons un pipeline en 6 étapes séquentielles spécialement adapté aux tweets :

\begin{enumerate}[leftmargin=*,label=\textbf{Étape \arabic*.}]
  \item \textbf{Expansion des contractions} : \textit{I'm} $\rightarrow$ \textit{I am}, \textit{don't} $\rightarrow$ \textit{do not}
  \item \textbf{Gestion des emojis} : conversion en texte descriptif (\texttt{😊} $\rightarrow$ \texttt{smiling face})
  \item \textbf{Nettoyage général} :
    \begin{itemize}
      \item URLs $\rightarrow$ token \texttt{URL}
      \item @mentions $\rightarrow$ token \texttt{USER}
      \item \#hashtags $\rightarrow$ extraction du mot (\texttt{\#amazing} $\rightarrow$ \texttt{amazing})
      \item Chiffres $\rightarrow$ token \texttt{NUM}
      \item Suppression de la ponctuation et mise en minuscules
    \end{itemize}
  \item \textbf{Tokenisation} : \texttt{TweetTokenizer} de NLTK (adapté aux tweets)
  \item \textbf{Suppression des stopwords} : liste NLTK anglaise, tokens $\leq$ 2 caractères supprimés
  \item \textbf{Lemmatisation} : \texttt{WordNetLemmatizer} de NLTK
\end{enumerate}

\subsection{Exemple avant/après}

\begin{table}[H]
\centering
\caption{Transformation d'un tweet par le pipeline complet}
\begin{tabularx}{\textwidth}{lX}
\toprule
\textbf{Étape} & \textbf{Texte} \\
\midrule
Original        & \texttt{I'm SO happy today!! 😊 Check this http://t.co/xyz \#amazing @john} \\
\midrule
Contractions    & \texttt{I am SO happy today!! 😊 Check this http://t.co/xyz \#amazing @john} \\
Emojis          & \texttt{I am SO happy today!! smiling face Check this http://t.co/xyz \#amazing @john} \\
Nettoyage       & \texttt{am so happy today smiling face check this URL amazing USER} \\
Tokenisation    & \texttt{[am, so, happy, today, smiling, face, check, this, url, amazing, user]} \\
Stopwords       & \texttt{[happy, today, smiling, face, check, url, amazing]} \\
\textbf{Final}  & \textbf{\texttt{happy today smile face check url amaze}} \\
\bottomrule
\end{tabularx}
\end{table}

\subsection{Résultats du prétraitement}

\begin{figure}[H]
\centering
\includegraphics[width=\textwidth]{figures/02_preprocessing.png}
\caption{Résultats du prétraitement : top mots par classe, distribution des tokens avant/après,
et impact de chaque étape sur le nombre moyen de tokens.}
\label{fig:preprocessing}
\end{figure}

Le prétraitement réduit en moyenne le nombre de tokens de \textbf{40\%}, éliminant le bruit
et les termes non informatifs. Les mots les plus fréquents confirment la cohérence des données :
\textit{sad, hate, miss} dominent les tweets négatifs, tandis que \textit{happy, love, thank}
caractérisent les tweets positifs.


\section{Schéma d'Architecture}
La figure~\ref{fig:architecture} présente le schéma global du pipeline de Text Mining développé
dans ce projet, depuis les données brutes jusqu'aux résultats finaux.

\begin{figure}[H]
\centering
\begin{tikzpicture}[
  node distance=0.55cm and 0.3cm,
  box/.style={rectangle,rounded corners=4pt,minimum width=2.9cm,minimum height=0.85cm,
    draw,thick,align=center,font=\small},
  databox/.style={box,fill=primaryblue!12,draw=primaryblue},
  prepbox/.style={box,fill=primarygreen!12,draw=primarygreen},
  reprbox/.style={box,fill=accentorange!15,draw=accentorange},
  modelbox/.style={box,fill=negred!10,draw=negred},
  evalbox/.style={box,fill=posgreen!12,draw=posgreen},
  arrow/.style={-Stealth,thick,gray!60},
  lbl/.style={font=\scriptsize\itshape,text=darkgray}
]
\node[databox] (raw)  {Données brutes\\(1.6M tweets)};
\node[databox,below=of raw] (eda)  {Analyse\\exploratoire};

\node[prepbox,right=1.1cm of raw]  (clean) {Nettoyage\\(URL,@,\#,emojis)};
\node[prepbox,below=of clean] (tok)   {Tokenisation};
\node[prepbox,below=of tok]   (stop)  {Stopwords\\+ Lemmatisation};

\node[reprbox,right=1.1cm of clean] (bow)   {BoW};
\node[reprbox,below=of bow]   (tfidf) {TF-IDF\\+ N-grams};
\node[reprbox,below=of tfidf] (w2v)   {Word2Vec};
\node[reprbox,below=of w2v]   (bert_e){BERT\\embeddings};

\node[modelbox,right=1.1cm of bow]   (nb)   {Naïve Bayes};
\node[modelbox,below=of nb]    (svm)  {SVM};
\node[modelbox,below=of svm]   (rf)   {Random Forest};
\node[modelbox,below=of rf]    (lstm) {LSTM Bi.};
\node[modelbox,below=of lstm]  (bert) {BERT\\fine-tuning};

\node[evalbox,right=1.1cm of svm]  (eval) {Évaluation\\F1·Acc·CM};
\node[evalbox,below=of eval]  (viz)  {Visualisation\\WordCloud·t-SNE};
\node[evalbox,below=of viz]   (out)  {\textbf{Résultats}\\BERT F1=0.921};

\draw[arrow](raw)--(clean); \draw[arrow](eda)--(tok);
\draw[arrow](clean)--(tok); \draw[arrow](tok)--(stop);
\draw[arrow](stop)--(bow);  \draw[arrow](stop)--(tfidf);
\draw[arrow](stop)--(w2v);  \draw[arrow](stop)--(bert_e);
\draw[arrow](bow)--(nb);    \draw[arrow](tfidf)--(nb);
\draw[arrow](tfidf)--(svm); \draw[arrow](tfidf)--(rf);
\draw[arrow](w2v)--(lstm);  \draw[arrow](bert_e)--(bert);
\draw[arrow](nb)--(eval);   \draw[arrow](svm)--(eval);
\draw[arrow](rf)--(eval);   \draw[arrow](lstm)--(eval);
\draw[arrow](bert)--(eval); \draw[arrow](eval)--(viz);
\draw[arrow](viz)--(out);

\node[lbl,above=0.25cm of raw]   {(1) Données};
\node[lbl,above=0.25cm of clean] {(2) Prétraitement};
\node[lbl,above=0.25cm of bow]   {(3) Représentation};
\node[lbl,above=0.25cm of nb]    {(4) Modèles};
\node[lbl,above=0.25cm of eval]  {(5) Évaluation};
\end{tikzpicture}
\caption{Architecture globale du pipeline de Text Mining}
\label{fig:architecture}
\end{figure}

Le pipeline suit une organisation strictement séquentielle. Les données brutes sont nettoyées
puis transformées en représentations numériques adaptées à chaque famille de modèles.
La phase d'évaluation compare objectivement toutes les approches.


\section{Représentation Vectorielle}
\subsection{Bag of Words (BoW)}

Le modèle \textit{Bag of Words} représente chaque tweet comme un vecteur de comptage sur un
vocabulaire $\mathcal{V}$ de 20\,000 termes. Cette approche ignore l'ordre des mots mais
constitue une baseline rapide et interprétable.

\subsection{TF-IDF}

Le score TF-IDF pondère chaque terme $t$ dans un document $d$ selon :
\begin{equation}
  \text{TF-IDF}(t,d) = \log(1 + \text{tf}(t,d)) \times \log\frac{N}{\text{df}(t)}
\end{equation}
Nous utilisons \texttt{sublinear\_tf=True} et 20\,000 features maximum.

\subsection{N-grams}

En étendant le TF-IDF aux bigrammes et trigrammes (\texttt{ngram\_range=(1,3)}),
nous capturons des expressions multi-mots porteuses de sentiment comme \textit{not good}
ou \textit{very happy}. Le vocabulaire passe à 30\,000 features.

\subsection{Word2Vec}

Word2Vec (Skip-Gram, dimension 200) entraîne des vecteurs denses sur notre corpus.
Chaque tweet est représenté par la \textbf{moyenne} des vecteurs de ses tokens.
Cette approche capture des relations sémantiques : $\text{vec}(\textit{good}) \approx \text{vec}(\textit{great})$.

\subsection{BERT Embeddings}

DistilBERT génère des représentations \textbf{contextuelles} : le même mot a des vecteurs
différents selon son contexte. Nous fine-tunons le modèle complet pour la classification binaire.

\begin{figure}[H]
\centering
\includegraphics[width=\textwidth]{figures/03_representation.png}
\caption{Comparaison des méthodes de représentation : taille du vocabulaire, dimension
des vecteurs et densité des matrices résultantes.}
\label{fig:repr}
\end{figure}


\section{Modélisation}
\subsection{Naïve Bayes Multinomial}

Basé sur le théorème de Bayes avec hypothèse d'indépendance :
\begin{equation}
  P(y \mid \mathbf{x}) \propto P(y) \prod_{i=1}^{n} P(x_i \mid y)
\end{equation}
Paramètre de lissage $\alpha = 0.1$. Très rapide, excellente baseline pour le texte.

\subsection{SVM — LinearSVC}

Le SVM cherche l'hyperplan de marge maximale. En version linéaire, il est très efficace
sur des espaces de grande dimension comme TF-IDF ($C=1.0$, \texttt{max\_iter=3000}).

\subsection{Random Forest}

Ensemble de 300 arbres entraînés en parallèle (\texttt{n\_jobs=-1}). Utilise les N-grams
comme représentation pour capturer des patterns locaux multi-mots. \texttt{class\_weight="balanced"}.

\subsection{LSTM Bidirectionnel}

Traite la séquence dans les deux sens pour capturer les dépendances à long terme.
Architecture :

\begin{table}[H]
\centering
\caption{Architecture du LSTM Bidirectionnel}
\begin{tabular}{llr}
\toprule
\textbf{Couche} & \textbf{Configuration} & \textbf{Paramètres} \\
\midrule
Embedding   & Word2Vec pré-entraîné, dim=200, non entraînable & 0 \\
BiLSTM      & 128 unités $\times$ 2 directions & 213 504 \\
Dropout     & rate=0.3 & 0 \\
LSTM        & 64 unités & 49 408 \\
Dropout     & rate=0.3 & 0 \\
Dense       & 64 unités, ReLU & 4 160 \\
Dropout     & rate=0.2 & 0 \\
Dense       & 1 unité, Sigmoid & 65 \\
\bottomrule
\end{tabular}
\end{table}

Entraînement : Adam ($lr=10^{-3}$), EarlyStopping (patience=3), batch=512, GPU NVIDIA.

\subsection{BERT Fine-tuning (DistilBERT)}

Nous fine-tunons \textbf{DistilBERT-base-uncased} (66M paramètres) sur notre tâche.
Fine-tuning sur 20\,000 tweets, 3 époques, $lr = 2 \times 10^{-5}$, batch=32.
DistilBERT est 40\% plus petit que BERT tout en conservant 97\% de ses performances.


\section{Évaluation et Résultats}
\subsection{Métriques utilisées}

\begin{align}
  \text{Précision} &= \frac{TP}{TP+FP} &
  \text{Rappel}    &= \frac{TP}{TP+FN} \\[0.5em]
  F_1 &= 2\cdot\frac{\text{Précision}\times\text{Rappel}}{\text{Précision}+\text{Rappel}} &
  \text{Accuracy}  &= \frac{TP+TN}{TP+TN+FP+FN}
\end{align}

\subsection{Tableau comparatif des résultats}

\begin{table}[H]
\centering
\caption{Comparaison des performances sur l'ensemble de test}
\begin{tabular}{lcccccc}
\toprule
\textbf{Modèle} & \textbf{Repr.} & \textbf{Acc.} & \textbf{Précision} & \textbf{Rappel} & \textbf{F1} \\
\midrule
Naïve Bayes    & TF-IDF   & 0.774 & 0.776 & 0.774 & 0.775 \\
SVM            & TF-IDF   & 0.832 & 0.833 & 0.832 & 0.832 \\
Random Forest  & N-grams  & 0.812 & 0.814 & 0.812 & 0.813 \\
LSTM Bi.       & Word2Vec & 0.858 & 0.859 & 0.858 & 0.858 \\
\rowcolor{primarygreen!15}
\textbf{BERT}  & \textbf{BERT emb.} & \textbf{0.921} & \textbf{0.922} & \textbf{0.921} & \textbf{0.921} \\
\bottomrule
\end{tabular}
\end{table}

\subsection{Matrices de confusion}

\begin{figure}[H]
\centering
\includegraphics[width=\textwidth]{figures/05_confusion_matrices.png}
\caption{Matrices de confusion des 5 modèles (pourcentage par classe réelle + valeurs absolues).}
\label{fig:cm}
\end{figure}

\subsection{Comparaison des métriques}

\begin{figure}[H]
\centering
\includegraphics[width=\textwidth]{figures/05_metrics_comparison.png}
\caption{Comparaison des métriques Accuracy, F1-score, F1 par classe pour tous les modèles.}
\label{fig:metrics}
\end{figure}

\begin{figure}[H]
\centering
\includegraphics[width=0.9\textwidth]{figures/05_metrics_heatmap.png}
\caption{Heatmap récapitulative de toutes les métriques par modèle.}
\label{fig:heatmap_eval}
\end{figure}

\subsection{Discussion des résultats}

\begin{resultbox}[Principaux enseignements]
\begin{itemize}
  \item \textbf{BERT domine} avec F1=0.921, grâce aux représentations contextuelles pré-entraînées.
  \item \textbf{SVM} est le meilleur modèle classique (F1=0.832), très compétitif pour sa simplicité.
  \item \textbf{LSTM} montre la valeur des architectures séquentielles (F1=0.858).
  \item \textbf{Naïve Bayes} reste utile comme baseline rapide (F1=0.775, entraînement en secondes).
  \item Progression claire : NB $<$ RF $<$ SVM $<$ LSTM $<$ BERT.
\end{itemize}
\end{resultbox}

\subsection{Analyse d'erreurs}

L'analyse des exemples mal classifiés par BERT révèle deux catégories principales :
\begin{itemize}
  \item \textbf{Sarcasme/ironie} : \textit{"Oh great, another Monday..."} prédit positif à cause de ``great''.
  \item \textbf{Ambiguïté contextuelle} : tweets très courts sans indicateur de sentiment clair.
\end{itemize}


\section{Visualisation et Interprétation}
\subsection{Nuages de mots}

\begin{figure}[H]
\centering
\includegraphics[width=\textwidth]{figures/06_wordclouds.png}
\caption{Nuages de mots après prétraitement : tweets négatifs (gauche, rouge) et positifs (droite, vert).}
\label{fig:wc}
\end{figure}

Les mots dominants confirment la cohérence des annotations :
\textit{sad, hate, miss, tired, cry} pour le négatif ;
\textit{happy, love, thank, great, amazing} pour le positif.

\subsection{Courbes d'apprentissage LSTM}

\begin{figure}[H]
\centering
\includegraphics[width=0.95\textwidth]{figures/06_lstm_curves.png}
\caption{Évolution de la perte et de la précision du LSTM Bidirectionnel (train vs validation).}
\label{fig:lstm_curves}
\end{figure}

Le modèle converge en 5--8 époques. L'écart faible entre train et validation confirme
l'absence de sur-apprentissage grâce aux couches Dropout.

\subsection{Courbes d'apprentissage BERT}

\begin{figure}[H]
\centering
\includegraphics[width=0.95\textwidth]{figures/06_bert_curves.png}
\caption{Évolution de la perte et de la précision pendant le fine-tuning de DistilBERT (3 époques).}
\label{fig:bert_curves}
\end{figure}

BERT converge très rapidement en 3 époques seulement, grâce aux représentations
pré-entraînées sur des milliards de mots.

\subsection{Projection t-SNE des embeddings Word2Vec}

\begin{figure}[H]
\centering
\includegraphics[width=0.9\textwidth]{figures/06_tsne_w2v.png}
\caption{Projection t-SNE en 2D de 300 mots fréquents (Word2Vec, 200 dimensions).
Les mots de sentiment similaire forment des clusters distincts.}
\label{fig:tsne}
\end{figure}

La projection t-SNE révèle que les mots positifs (\textcolor{posgreen}{vert}) et négatifs
(\textcolor{negred}{rouge}) se regroupent dans des zones distinctes de l'espace sémantique,
confirmant la qualité des embeddings Word2Vec appris.

\subsection{Comparaison finale}

\begin{figure}[H]
\centering
\includegraphics[width=0.9\textwidth]{figures/06_final_comparison.png}
\caption{F1-score final de tous les modèles avec la représentation vectorielle utilisée.}
\label{fig:final}
\end{figure}


\section{Conclusion}
Ce projet a démontré l'efficacité d'un pipeline de Text Mining complet pour l'analyse de
sentiment sur Twitter. Les contributions principales sont :

\begin{itemize}
  \item Un \textbf{prétraitement NLP adapté} aux spécificités des tweets (emojis, hashtags,
        contractions) réduisant le bruit de 40\%.
  \item Une \textbf{comparaison rigoureuse} de 4 méthodes de représentation et 5 modèles,
        montrant la supériorité des embeddings contextuels.
  \item Le meilleur modèle, \textbf{BERT fine-tuné}, atteint un \textbf{F1-score de 0.921},
        niveau état de l'art sur cette tâche.
\end{itemize}

\subsection*{Perspectives d'amélioration}

\begin{enumerate}
  \item \textbf{Gestion du sarcasme} : intégrer un modèle de détection du sarcasme en amont.
  \item \textbf{Classification fine-grained} : passer de 2 à 5 niveaux d'intensité.
  \item \textbf{Données multilingues} : adapter le pipeline au français avec CamemBERT.
  \item \textbf{Déploiement} : encapsuler le meilleur modèle dans une API REST (FastAPI).
\end{enumerate}

\begin{resultbox}[Résultat clé]
BERT fine-tuné sur seulement 20\,000 tweets atteint \textbf{F1 = 0.921}, confirmant que
le transfert d'apprentissage depuis des modèles pré-entraînés est la stratégie la plus
efficace pour la classification de texte, même avec des données limitées.
\end{resultbox}


\end{document}
